\begin{theorem}[Source-fidelity census]\label{thm:source-fidelity-census}
For the frozen HodgeCY v1.0.0 double-octic cohort, the HodgeCY II source
census processes \(456\) presentations and records \(114\) nontrivial
source-level pairs or larger sets.  These decompose as \(57\) pairs,
\(13\) triples, and \(44\) larger sets.
\end{theorem}

\begin{proof}[Evidence]
This is a frozen computational census statement.  The reproducible evidence is
the final HodgeCY II synthesis bundle and the deterministic reproduction script
`scripts/reproduce_hodgecy_ii.py`.
\end{proof}

\begin{proposition}[The \(84/84a\) source witness]\label{prop:source-witness-84-84a}
The presentations \(84\) and \(84a\) have the same local inventory, the same
Hodge data \((h^{1,1},h^{1,2})=(40,0)\), and the same rational source rank
\(26\).  They are separated by integral source type:
\[
  {\rm SNF}(84)=(1^{23},2,6,12),
  \qquad
  {\rm SNF}(84a)=(1^{21},2,4,4,4,12).
\]
\end{proposition}

\begin{proof}[Evidence]
This is the source-lattice comparison certified in the HodgeCY II final
evidence bundle.  The proposition is a statement about the verified source
assembly records only.
\end{proof}

\begin{theorem}[Eight-plane line skeleton]\label{thm:line-skeleton}
Let \(S=k[x_0,x_1,x_2,x_3]\), and let
\(L_1,\ldots,L_8\in S_1\) be distinct linear forms such that every pair has
rank \(2\) and every triple has rank \(3\).  Equivalently, no three
corresponding planes contain a common line.  If \(F=\prod_i L_i\) and
\[
  I_C=(F/L_1,\ldots,F/L_8),
\]
then
\[
  I_C=\bigcap_{1\le i<j\le 8}(L_i,L_j).
\]
Thus \(I_C\) is the saturated reduced ideal of the union of the \(28\)
pairwise intersection lines, and it has Hilbert--Burch resolution
\[
0\longrightarrow S(-8)^7\longrightarrow S(-7)^8
\longrightarrow I_C\longrightarrow 0 .
\]
Four-plane, or higher, point concurrences are allowed: the argument only uses
the codimension-two condition that no three planes contain a common line.
\end{theorem}

\begin{proof}[Proof source]
Use the self-contained induction in
`research_outputs/hodgecy_ii/final/hodgecy_ii_hilbert_burch_theorem.tex`.
For the manuscript body, reproduce that proof rather than replacing it by an
appeal to the full star-configuration theorem of
Geramita--Harbourne--Migliore~\cite{GeramitaHarbourneMigliore2013}.
\end{proof}

\begin{proposition}[Regular quartic block section and mapping cone]\label{prop:quartic-block-section}
For the verified HodgeCY II arrangements \(84\) and \(84a\), the frozen quartic
block section is regular on the eight-plane line skeleton.  If \(I_B=I_C+(Q)\),
then multiplication by \(Q\) gives an exact sequence
\[
0\longrightarrow (S/I_C)(-4)\longrightarrow S/I_C
\longrightarrow S/I_B\longrightarrow 0,
\]
and \(S/I_B\) has mapping-cone resolution
\[
0
\longrightarrow S(-12)^7
\longrightarrow S(-11)^8 \oplus S(-8)^7
\longrightarrow S(-7)^8 \oplus S(-4)
\longrightarrow S
\longrightarrow S/I_B
\longrightarrow 0 .
\]
\end{proposition}

\begin{proof}[Proof source]
This is the regular quartic block-section theorem and mapping-cone corollary
in the final Hilbert--Burch theorem asset.  The verifier checks nonzero
squarefree restriction of \(Q\) on all \(28\) associated line primes for both
frozen arrangements.
\end{proof}

\begin{corollary}[Verified block Hilbert series]\label{cor:block-hilbert-series}
For the verified HodgeCY II block schemes attached to \(84\) and \(84a\),
\[
  {\rm Hilb}_{S/I_B}(t)=
  \frac{(1-t^4)(1-8t^7+7t^8)}{(1-t)^4}.
\]
Consequently
\[
H_B(d)=1,4,10,20,34,52,74,92,105,112,112,\ldots,
\]
with stabilization at degree \(9\).
\end{corollary}

\begin{proof}[Proof source]
This follows from the verified regular quartic block section and the mapping
cone argument in the final Hilbert--Burch theorem asset.
\end{proof}

\begin{corollary}[Degree-eight block-evaluation deficiency]\label{cor:block-deficiency}
For the verified HodgeCY II block schemes attached to \(84\) and \(84a\),
\[
  h^1({\bf P}^3,\mathcal I_B(8))=7.
\]
Equivalently, the degree-eight block-evaluation cokernel has dimension \(7\).
This is a block-scheme deficiency statement only.
\end{corollary}

\begin{proposition}[Source versus block-evaluation non-determination]\label{prop:source-block-nondetermination}
On the witness pair \(84/84a\), the verified degree-eight block-evaluation data
do not determine the integral source-assembly type.  The source layer splits by
Smith type, while the verified block Hilbert/evaluation data agree through
degree \(8\).
\end{proposition}

\begin{proof}[Evidence]
Combine Proposition~\ref{prop:source-witness-84-84a} with
Corollaries~\ref{cor:block-hilbert-series} and~\ref{cor:block-deficiency}.
This does not assert a reverse implication, a canonical source-to-evaluation
map, an integral evaluation lattice, an LMHS identification, or a classical
nodal-defect theorem.
\end{proof}
